\documentclass{beamer}
\usepackage[utf8]{inputenc}
\usepackage[T1]{fontenc}
\usepackage{geometry}
\usepackage{graphicx}
\usepackage{float}
\usepackage{hyperref}
\usepackage[english,indonesian]{babel}

\title{Optimasi Pemrosesan Template LaTeX Menggunakan Pendekatan Berbasis AI}
\author{Academic Support Agent}
\date{\today}

\begin{document}
\maketitle

\begin{abstract}

\end{abstract}

\chapter{PENDAHULUAN}


\chapter{TINJAUAN PUSTAKA}


% Original content (commented out):
% % Created 2017-12-14 jeu. 14:25
% Intended LaTeX compiler: pdflatex
\documentclass[10pt,svgnames,fragile]{beamer}
                \usepackage[english]{babel}
\usepackage[utf8x]{inputenc}
\usepackage{lmodern}
\usepackage[T1]{fontenc}
\usepackage{graphicx}
\usepackage{etex}
\usepackage{xcolor}
\usepackage[normalem]{ulem}
\usepackage{textcomp}
\usepackage{pdflscape}
\usepackage{marvosym}
\usepackage{wasysym}
\usepackage{amssymb}
\usepackage{amsmath}
\usepackage{amsthm}
\usepackage{qtree}
\usepackage{bussproofs}
\usepackage{proof}
\usepackage{fitch}
\usepackage{cancel}
\usepackage{url}
\usepackage{smfthm}
\AtBeginSection[]{\begin{frame}<beamer>\frametitle{}\tableofcontents[currentsection,hideothersubsections]\end{frame}}
\subtitle{}
\institute[Université de Lorraine]{Département de philosophie \\ Archives Henri Poincaré - UMR 7117 du CNRS \\ Université de Lorraine \\ 91 bd Libération, 54000 Nancy \\ France }
\titlegraphic{\includegraphics[height=1cm]{ahp}\includegraphics[height=1cm]{udl}\includegraphics[height=1cm]{cnrs}}
\usetheme{CambridgeUS}
\usepackage{beamer_udl_theme}
\setbeamertemplate{navigation symbols}{%
	%insertslidenavigationsymbol%
	%insertframenavigationsymbol%
	%insertsubsectionnavigationsymbol%
	%insertsectionnavigationsymbol%
	%insertdocnavigationsymbol%
	%insertbackfindforwardnavigationsymbol%
}
\usetheme{default}
\author{Joseph Vidal-Rosset}
\date{}
\title{Presentation template}
\begin{document}

\maketitle

\section{Section 1}
\label{sec:org481dcf5}
\begin{frame}[label={sec:org9a582a5}]{}
First point 
\end{frame}

\section{Section 2}
\label{sec:org62f0986}
\begin{frame}[label={sec:orgea71ea0}]{}
\begin{block}{A title}
Text and text again \ldots{} 
\end{block}
\end{frame}
\section{Section 3}
\label{sec:org69f6999}
\begin{frame}[label={sec:orgb86c640}]{A bigger title}
Text\ldots{}  \pause 
\begin{block}{A subtitle}
Text and bibliographical reference, for example about  Quine \cite{quine1934}
\end{block}
\end{frame}

\begin{frame}[label={sec:org3031881}]{References}
\bibliography{reforg}
\bibliographystyle{smfplain}
\end{frame}
\end{document}

\end{document}